% Propietario conservado: /Users/ruben/Documents/New project/output/REVISION_ARGUMENTAL_APERTURAS_CIERRES_20260915/VI/delivery/MANUSCRIPT_VI_ES_EN_COMPLETE/source_es/sections/07_contraste_metrologico.tex
% Artículo VI. Propietario: indice_articulos. Adición autónoma de exposición.
% Procedencia: RESULTADO_RECUPERADO; controles documentales nuevos separados.
% REV2/fuente/modulos/15_comparacion_cuantitativa_integral.md (completo).
% REV2/fuente/modulos/16_precision_robustez_y_contraste.md (completo).
% REV2/fuente/modulos/19_realizacion_exhaustiva_estado_por_estado.md (completo).
% Integral: integracion_83/masas_p0/corpus/source/masas/
% 19_comparacion_cuantitativa_posterior.tex, parte anterior a \endinput.
% Integral: integracion_83/masas_p0/tex/apendice_inventario_471_384_sucesor.tex.
% Depende de la construcción de atlas, firma, torres y realización material
% incorporada por el editor en las secciones precedentes. No la sustituye.

\section{Realización de los observables y contraste metrológico}
\label{vi:sec:vi-contraste-metrologico}

El valor de una masa adquiere su significado dentro de la estructura que lo
produce y de la representación en la que se mide. La construcción
APP--TRIT--TPK determina primero las trayectorias admisibles, sus registros
orientados y los lectores de la firma; las torres y la realización material
actúan sobre esos objetos. El contraste experimental se sitúa al término de
esta composición. Su función consiste en comparar observables del mismo
tipo, conservar la identidad del estado y cuantificar el acuerdo en una
carta común. Esta posición permite estudiar conjuntamente la estructura
matemática, la coordenada evaluada y su interpretación física.

La presente sección reúne tres niveles que cumplen funciones diferentes:
el dominio interno de las rutas y sus cocientes; las clases y evaluaciones
individualizadas en los propietarios de masas; y el inventario externo de
observables utilizado como referencia. El orden de exposición conserva la
dirección constructiva del artículo. Las magnitudes externas comparecen como
términos de contraste, mientras los valores internos conservan las firmas,
operaciones y condiciones que fijan su evaluación.

\subsection{Dominio estructural, clases materiales e inventario externo}
\label{vi:subsec:vi-dominios-contraste}

Sean \(\mathcal C\) el conjunto de celdas APP y \(\mathcal R\) su dominio
de rutas orientadas. La identificación de una ruta contiene las coordenadas
de celda y las dos orientaciones iniciales:
\[
 \gamma=(r,c,h_0,v_0),\qquad
 r,c\in\{1,\ldots,9\},\quad h_0,v_0\in\{-1,+1\}.
\]
Por tanto, \(\#\mathcal C=81\) y \(\#\mathcal R=81\cdot4=324\).
Las \(56\) familias de trayectoria y las \(13\) multisecciones pertenecen
a cocientes y agrupaciones estructurales del mismo dominio, con las
relaciones y lectores establecidos anteriormente. Estos cuatro cardinales
describen objetos internos distintos. La asignación de una representación
física puede conservar una fibra completa; una etiqueta de especie no
impone por sí misma una elección puntual dentro de ella.

La tabla pública de \(23\) clases organiza las realizaciones por sector.
Su función es relacionar la residencia interna con la representación, la
carga, el espín, el color, el sabor y el tipo de observable. Una clase puede
comprender varias especies, y una identidad catalográfica puede reunir
distintos estados de carga. La correspondencia se expresa mediante mapas
tipados y conserva estas multiplicidades.

El inventario metrológico congelado contiene \(471\) registros de masa y
\(384\) de anchura. El corpus documenta \(1170\) identidades en \(438\)
grupos, de los cuales \(403\) contienen al menos uno de esos observables.
La suma
\[
 471+384=855
\]
cuenta observables externos de tabla-resumen. El atlas \(81/56/13/324\),
las \(23\) clases y los \(855\) registros pertenecen así a tres niveles
de descripción. Su asociación conserva el tipo de cada conjunto, en lugar
de convertir sus cardinales en un único número de partículas o predicciones.

\begin{center}
\begin{tabular}{p{0.49\linewidth}rrr}
\toprule
Familia catalográfica & Grupos & Masas & Anchuras\\
\midrule
General & 46 & 53 & 9\\
Quarks & 6 & 8 & 1\\
Mesones & 201 & 243 & 225\\
Bariones & 149 & 166 & 149\\
Registro gravitatorio externo & 1 & 1 & 0\\
\midrule
Total con observable & 403 & 471 & 384\\
\bottomrule
\end{tabular}
\end{center}

La fila gravitatoria conserva la identidad que figura en la fuente externa;
su presencia en el inventario no define una partícula elemental de HMT.
Los \(35\) grupos restantes conservan su identidad catalográfica, aunque
el corte no les asigne una masa o anchura de tabla-resumen. La integridad
documental abarca ambos hechos: las magnitudes que constan y las identidades
para las que ese corte no publica una de dichas magnitudes.

\subsection{El observable antes de la coordenada numérica}
\label{vi:subsec:vi-tipos-observable}

Una salida material se especifica mediante el estado, la representación,
el operador y el procedimiento de lectura. Para una fibra \(F_X\), sea
\(\widehat M_X\) el operador de masa ya construido y sea \(P_{X,a}\)
el proyector correspondiente al estado o canal de preparación \(a\).
Para \(\Psi_X\) normalizado en su imagen, la distribución espectral es
\[
 \mu_{X,a}(B)
 =\langle\Psi_X,
   P_{X,a}E_{\widehat M_X}(B)P_{X,a}\Psi_X\rangle,
\]
donde \(E_{\widehat M_X}\) es la medida espectral del operador. Esta
expresión conserva el estado preparado y la multiplicidad de la fibra.
Una coordenada escalar única para todo estado preparado en ese subespacio
se obtiene bajo la condición
\[
 \operatorname{Ran}(P_{X,a})\subset\operatorname{Dom}(\widehat M_X),
 \qquad \widehat M_XP_{X,a}=m_XP_{X,a}.
\]
En efecto, cada vector de ese subespacio es entonces un vector propio del
operador de masa para el mismo valor \(m_X\), y su medida espectral está
concentrada en ese valor. Si se emplea la igualdad de compresión
\(P_{X,a}\widehat M_XP_{X,a}=m_XP_{X,a}\), se exige además que
\(P_{X,a}\) reduzca \(\widehat M_X\); en dimensión finita, equivale
a su conmutación. Sin esa condición, la compresión puede fijar únicamente
el valor medio. Cuando la medida espectral mantiene varios valores, la
salida se presenta como distribución o multiplete con sus pesos.

Para una densidad fija \(\rho\) en una fibra finita, sean \(m_j\)
los valores propios distintos y \(E_j\) sus proyectores espectrales.
Los pesos \(p_j=\operatorname{tr}(\rho E_j)\geq0\), con
\(\sum_jp_j=1\), determinan
\[
 \overline m=\sum_jp_jm_j,\qquad
 \operatorname{Var}_{\mu}(m)=\sum_jp_j(m_j-\overline m)^2.
\]
\paragraph{Criterio de concentración escalar para el estado fijado.}
La varianza se anula si y sólo si todos los valores con peso positivo
coinciden. En efecto, cada sumando es no negativo; la suma es cero
precisamente cuando \(m_j=\overline m\) para todo \(j\) con
\(p_j>0\). Así queda determinado cuándo una distribución de masa admite
una publicación escalar sin pérdida de la dispersión del estado preparado.
Este criterio se refiere a la densidad fijada y se distingue de la condición
operatorial sobre todo el subespacio. El caso electrónico utiliza su
operador sobre la fibra central, tal como se ha construido previamente;
el rango uno de un proyector arbitrario no basta para concentrar la medida
espectral de otro operador.

Masa y anchura utilizan lectores coordinados de tipos diferentes. Para un
estado resonante, una convención de polo puede expresarse en coordenadas de
energía como
\[
 z_X=M_Xc_{\rm int}^{2}-\frac{i}{2}\Gamma_X.
\]
La parte real es la coordenada de masa-energía; \(\Gamma_X\) tiene
dimensión de energía y mide la anchura. Una vez fijada la carta temporal,
\(\tau_X=\hbar_{\rm int}/\Gamma_X\). La relación supone el régimen de
polo y supervivencia para el que ha sido construida; una tasa, una anchura y
un tiempo de vida mantienen sus unidades y sus reglas de conversión. Esta
distinción explica por qué el inventario conserva dos clases de observable.

El mismo principio determina la comparación por sector. Las masas corrientes
de quark se expresan con esquema \(\mathscr S\) y escala \(\mu\);
las resonancias distinguen masa de polo, parametrización Breit--Wigner y
convención sobre capa. Los neutrinos distinguen autovalores de masa,
diferencias de cuadrados, sabores y parámetros de mezcla. Las cotas
unilaterales conservan el nivel de confianza y la convención estadística
documentados. Una cota superior no se convierte en un valor central positivo,
ni una diferencia de cuadrados en una masa absoluta.

\subsection{Fijación de la construcción y correspondencia de registros}
\label{vi:subsec:vi-congelacion}

La realización de una especie \(X\) transmite la información necesaria
desde sus antecedentes internos. Podemos reunirla, a efectos de contraste,
en el registro
\[
 \mathfrak D_X=
 \bigl(F_X,\rho_X,P_{X,a},\mathcal L_X,
       \sigma_X,\eta_X,\widehat M_X,\mathcal U_X\bigr),
\]
donde \(\rho_X\) identifica la representación, \(\mathcal L_X\) el
registro genealógico, \(\sigma_X\in\mathbb Z^5\) la firma,
\(\eta_X\) la sección de torre y \(\mathcal U_X\) la carta de
unidades. Todos estos datos reciben su sentido de las operaciones
APP--TRIT--TPK y de las realizaciones demostradas en el artículo. Reunirlos
en una tupla es una operación de documentación: sus construcciones y
pruebas permanecen en las secciones que los producen.

La referencia externa se representa separadamente por
\[
 o=(\mathrm{id},\mathrm{especie},\mathrm{observable},
      v,\mathrm{unidad},\mathscr S,\mu,
      \mathrm{incertidumbre},\text{límite},\mathrm{confianza},
      \mathrm{fuente},\text{edición}).
\]
La asociación \((\mathfrak D_X,o)\) debe establecerse mediante la
identidad física y el tipo de observable. En una tabla extensa, esta
operación se realiza por claves de identidad y procedencia; la proximidad
de dos cifras no desempeña la función de clave. El registro mantiene la
etiqueta de partícula o antipartícula, el grupo que posee el observable y
la identidad de la propiedad medida.

La independencia del procedimiento de selección se examina ocultando o
permutando la columna metrológica y comprobando que \(\mathfrak D_X\)
permanece inalterado. Ese ensayo tiene sentido para un generador cuya
composición efectiva se puede ejecutar; la mera conservación de un archivo
de salida no constituye por sí misma tal ensayo. La huella criptográfica
fija qué artefacto se compara y hace detectables las modificaciones, pero
la independencia causal pertenece a la operación y a sus dependencias.

Los estatutos se asignan al dominio y a la operación que fundamentan cada
afirmación. Una fila puede documentar simultáneamente un operador interno,
una evaluación dada una firma y una comparación externa, sin identificar
entre sí esos tres niveles. Toda realización ulterior se vincula con la
definición, proposición o tabla que establece su alcance. En particular, las
clasificaciones documentales del catálogo se aplican al corte que especifica
cada registro y no modifican el estatuto matemático de una construcción
posterior.

\subsection{Correspondencia metrológica de las veintitrés clases}
\label{vi:subsec:vi-clases-metrologia}

La tabla siguiente conserva las \(23\) clases de la fuente y especifica la
forma de lectura que cada una exige al alcanzar la interfaz metrológica. Su
función difiere de la clasificación operatoria de la
proposición~\ref{vi:vi:prop:producto-fibrado-especie-ruta} y del inventario
documental de las tablas~\ref{vi:tab:vi-estados-clases} y~\ref{vi:tab:vi-clases}:
aquí se distinguen la magnitud contrastable, la convención de realización y
los datos que han de acompañarla. Las denominaciones de sabor y sector
orientan la correspondencia; el objeto matemático transmitido permanece en
las fibras, representaciones y lectores correspondientes.

\begingroup
\small
\setlength{\tabcolsep}{3pt}
\begin{longtable}{p{0.20\linewidth}p{0.25\linewidth}p{0.47\linewidth}}
\caption{Lectura metrológica de las veintitrés clases y tipo de objeto transmitido.}
\label{vi:tab:vi-clases-metrologia}\\
\toprule
Clase & Representantes & Objeto documentado para la lectura\\
\midrule
\endfirsthead
\toprule
Clase & Representantes & Objeto documentado para la lectura\\
\midrule
\endhead
\bottomrule
\endfoot
Leptón cargado 1 & Electrón y positrón & Ruta central, conjugación y coordenada beta.\\
Leptón cargado 2 & Muón y antimuón & Fibra, lectores y evaluación de la firma declarada.\\
Leptón cargado 3 & Tau y antitau & Fibra, lectores y evaluación de la firma declarada.\\
Neutrino 1 & Sabor electrónico y conjugado & Fibra neutra y realización masa--sabor.\\
Neutrino 2 & Sabor muónico y conjugado & Fibra neutra y realización masa--sabor.\\
Neutrino 3 & Sabor tauónico y conjugado & Fibra neutra y realización masa--sabor.\\
Quark arriba 1 & \(u,\bar u\) & Fibra de sabor y color; carta de esquema y escala.\\
Quark abajo 1 & \(d,\bar d\) & Fibra de sabor y color; carta de esquema y escala.\\
Quark arriba 2 & \(c,\bar c\) & Fibra de sabor y color; carta de esquema y escala.\\
Quark abajo 2 & \(s,\bar s\) & Fibra de sabor y color; carta de esquema y escala.\\
Quark arriba 3 & \(t,\bar t\) & Fibra de sabor y color; carta de esquema y escala.\\
Quark abajo 3 & \(b,\bar b\) & Fibra de sabor y color; carta de esquema y escala.\\
Calibre electromagnético & Fotón & Carta nula y representación de helicidad.\\
Calibre cromático & Gluones & Carta nula y representación adjunta de color.\\
Calibre cargado & \(W^+,W^-\) & Ruta, firma y convención de polo de la realización.\\
Calibre neutro & \(Z^0\) & Ruta, mezcla neutra y convención de polo.\\
Escalar & Higgs & Representación escalar, firma y convención de polo.\\
Compuesto mesónico & Mesones y conjugados & Constituyentes, registro de enlace y lectura compuesta.\\
Compuesto bariónico & Bariones y conjugados & Composición trilineal, enlace y lectura compuesta.\\
Topológico Fibonacci & \(1,\tau\) & Anillo de fusión y dimensión de Frobenius--Perron.\\
Topológico Ising & \(1,\psi,\sigma\) & Reglas de fusión y representación sectorial.\\
Topológico cíclico & \(\mathbb Z/6\mathbb Z\) & Caracteres de devanado y representación de trenzas.\\
Virtual & Secciones intermedias & Horizonte finito, prolongación y obstrucción.\\
\end{longtable}
\endgroup

En los leptones cargados, la lectura conduce a la masa de la sección central
o a la evaluación de la fibra seleccionada. En los neutrinos conserva además
la transformación entre masa y sabor y las diferencias cuadráticas; por ello
no se reduce a un escalar aislado. En los quarks, toda cifra se acompaña del
esquema y de la escala de renormalización. Las clases electromagnética y
cromática publican respectivamente una cota experimental y la sección nula
de calibre; las clases \(W\), \(Z\) y escalar requieren la convención de polo
o de resonancia declarada. Los compuestos conservan masa, anchura e identidad
de constituyentes. Las tres clases topológicas se contrastan mediante sus
leyes de fusión, trenzado y brecha una vez fijada la realización dinámica;
sus dimensiones y caracteres pertenecen a codominios propios. La clase
virtual conserva el polo, la obstrucción y el horizonte de prolongación. Esta
lectura determina qué objeto alcanza la metrología sin confundir las
veintitrés clases con veintitrés escalares universales de masa.

\subsection{Corte metrológico, unidades y exactitud de la transcripción}
\label{vi:subsec:vi-corte-metrologico}

Se utiliza el corte identificado por el corpus como Particle Data Group
2026, con fecha de corte \(15\) de enero de \(2026\) y consulta
documentada de \(26\) de julio de \(2026\). Esta sección reproduce ese
estado documental: no incorpora una actualización posterior del catálogo.
Las referencias del suplemento de evaluación y las de la malla histórica
mantienen su procedencia individual; no reciben automáticamente la edición
del inventario general.

Cada valor se conserva como cadena decimal original, junto con su
presentación, unidad, incertidumbre y clave de procedencia. Esta decisión
evita que una conversión intermedia a coma flotante cambie los últimos
dígitos o elimine ceros significativos de la transcripción. La cadena
original y la representación tipográfica son dos datos distintos: la
primera permite reproducir exactamente la aritmética del archivo; la
segunda comunica el resultado con una precisión adecuada para la lectura.
Los dígitos de representación digital de una referencia no aumentan su
precisión experimental.

La conversión entre MeV y GeV utiliza el factor exacto \(10^3\), aplicado
al valor y a su incertidumbre. Para presentar una masa, la unidad se escribe
MeV/\(c^2\) o GeV/\(c^2\); para presentar la energía de reposo,
MeV o GeV. La relación \(E=mc^2\) es la carta que permite pasar de una
lectura a la otra. Una comparación de coordenadas exige fijar esa carta en
ambas columnas.

El corte del inventario mantiene expresamente como no suministrados el
esquema y la escala de sus registros externos. El nombre de una propiedad
puede contener una indicación como ``Breit--Wigner''; se conserva literalmente,
pero esa indicación parcial no completa los metadatos que la comparación
requiere. Las casillas vacías, las marcas de dato no disponible y las de
no aplicabilidad se mantienen diferenciadas. Una incógnita documental
no se sustituye por cero.

\subsection{Evaluaciones individualizadas y condiciones de lectura}
\label{vi:subsec:vi-evaluaciones-documentadas}

La tabla~\ref{vi:tab:vi-evaluaciones-veinte} reúne \(20\) presentaciones
numéricas: la coordenada beta electrónica, dos salidas nulas y
\(17\) evaluaciones condicionadas a descriptores declarados. Dentro de
estas últimas, \(11\) corresponden a sectores elementales y \(6\) a
compuestos. Esta cuenta expresa el contenido de esa presentación documental;
las fibras, familias y registros del dominio conservan su extensión propia.

La tabla mantiene las veinte entradas. Para facilitar su lectura se
identifican cuatro modalidades: \(\mathrm{B}\), coordenada beta
electrónica; \(\mathrm{N}\), carta nula; \(\mathrm{F}\), evaluación
condicionada a una firma completa; \(\mathrm{A}\), evaluación angular
condicionada. El subíndice \(c\) señala la composición material. La
condición de cada entrada pertenece a su tipo de resultado y se conserva
también en el catálogo electrónico que acompaña el texto.

\begingroup
\small
\setlength{\tabcolsep}{3pt}
\begin{longtable}{p{0.13\linewidth}p{0.32\linewidth}p{0.38\linewidth}p{0.07\linewidth}}
\caption{Veinte evaluaciones individualizadas, con su carta y modalidad de contraste.}
\label{vi:tab:vi-evaluaciones-veinte}\\
\toprule
Estado & Coordenada HMT, MeV/\(c^2\) & Referencia documental posterior & Tipo\\
\midrule
\endfirsthead
\toprule
Estado & Coordenada HMT, MeV/\(c^2\) & Referencia documental posterior & Tipo\\
\midrule
\endhead
\bottomrule
\endfoot
\(e^\pm\) & \(0.510998950\allowbreak690000809\) & \(0.51099895069(16)\) MeV/\(c^2\) & B\\
\(\mu^\pm\) & \(105.658360294\allowbreak435829303\) & \(105.6583755\pm0.0000023\) MeV/\(c^2\) & F\\
\(\tau^\pm\) & \(1776.860917631\allowbreak432295595\) & \(1776.93\pm0.09\) MeV/\(c^2\) & F\\
\(u,\bar u\) & \(2.165924536\allowbreak173907663\) & \(2.16\pm0.07\) MeV/\(c^2\) & A\\
\(d,\bar d\) & \(4.679550162\allowbreak948874172\) & \(4.70\pm0.07\) MeV/\(c^2\) & A\\
\(c,\bar c\) & \(1270.500838641\allowbreak911135286\) & \(1.2729\pm0.0045\) GeV/\(c^2\) & A\\
\(s,\bar s\) & \(92.940093907\allowbreak845640641\) & \(92.9\pm0.7\) MeV/\(c^2\) & A\\
\(t,\bar t\) & \(172597.479544019\allowbreak136261367\) & \(172.60\pm0.27\) GeV/\(c^2\) & A\\
\(b,\bar b\) & \(4190.119763299\allowbreak790218075\) & \(4.186\pm0.006\) GeV/\(c^2\) & A\\
\(\gamma\) & \(0\) & \(<10^{-18}\) eV/\(c^2\), cota conservada & N\\
\(g\) & \(0\) & Carta de calibre nula; sin fila experimental de contraste & N\\
\(W^\pm\) & \(80378.964909704\allowbreak547024865\) & \(80.3625\pm0.0077\) GeV/\(c^2\) & F\\
\(Z^0\) & \(91187.533444313\allowbreak407844855\) & \(91.1879\pm0.0020\) GeV/\(c^2\) & F\\
\(H\) & \(125292.466042536\allowbreak133000433\) & \(125.13\pm0.11\) GeV/\(c^2\) & A\\
\(\pi^0\) & \(134.976724327\allowbreak872125739\) & \(134.976827767685\) MeV/\(c^2\), suplemento & \(\mathrm F_c\)\\
\(\pi^\pm\) & \(139.570485171\allowbreak572191375\) & \(139.570390983681\) MeV/\(c^2\), suplemento & \(\mathrm F_c\)\\
\(K^\pm\) & \(493.677085649\allowbreak530612476\) & \(493.676599458041\) MeV/\(c^2\), suplemento & \(\mathrm F_c\)\\
\(K^0\) & \(497.611186497\allowbreak750457359\) & \(497.610767531258\) MeV/\(c^2\), suplemento & \(\mathrm F_c\)\\
\(p,\bar p\) & \(938.271751546\allowbreak984898299\) & \(938.27208943\) MeV/\(c^2\), suplemento & \(\mathrm F_c\)\\
\(n,\bar n\) & \(939.565810472\allowbreak507023313\) & \(939.56542194\) MeV/\(c^2\), suplemento & \(\mathrm F_c\)\\
\end{longtable}
\endgroup

En esta tabla, los valores de referencia se presentan con la notación
decimal y las unidades declaradas para cada observable. El registro de
cálculo conserva además la cadena decimal empleada en la evaluación,
incluida la diferencia entre esa cadena y su presentación redondeada para
el tau, \(Z\) y Higgs. La comparación de los quarks
permanece informativa mientras no se reúna la carta de esquema y escala.
Los sectores \(W,Z,H\) requieren conservar también la convención de
resonancia. Las seis referencias compuestas proceden del suplemento
identificado en los registros de evaluación; sus metadatos no se completan
por analogía con las referencias elementales.

Las firmas completas documentadas en las modalidades \(\mathrm F\) y
\(\mathrm F_c\), en el orden
\(\sigma=(n_A,s_C,k_{\Delta_4},b_{120},b_\zeta)\), son:
\[
\begin{array}{c|r@{,\,}r@{,\,}r@{,\,}r@{,\,}r}
\mu&42&-3&8&0&2\\
\tau&64&0&25&-1&6\\
W&94&-1&0&-2&4\\
Z&95&-1&-11&0&-4\\
\pi^0&44&-4&-25&1&-2\\
\pi^\pm&44&1&-2&2&-1\\
K^\pm&54&-1&17&-2&2\\
K^0&54&0&32&3&-2\\
p&59&0&9&1&0\\
n&59&1&-34&0&1
\end{array}
\]
Las modalidades angulares conservan, respectivamente para
\(u,d,c,s,t,b,H\), los pares
\[
(n_A,s_C)=(11,7),(17,8),(61,8),(41,-3),(100,-1),(71,-5),(97,9).
\]
Los dos componentes declarados no se amplían con
tres ceros implícitos. La ecuación de evaluación y el mecanismo que
produce un descriptor son afirmaciones diferentes; las pruebas de ambas
se encuentran en sus definiciones y proposiciones respectivas.

\subsection{Diferencias relativas y propagación de incertidumbre}
\label{vi:subsec:vi-residuales}

Fijada una carta común y dos coordenadas positivas del mismo observable,
la diferencia relativa y la diferencia logarítmica son
\[
 \delta_X=\frac{m_X^{\rm HMT}-m_X^{\rm ext}}{m_X^{\rm ext}},
 \qquad r_X=\log\frac{m_X^{\rm HMT}}{m_X^{\rm ext}},
 \qquad \delta_X^{\rm ppm}=10^6\delta_X.
\]
Una conversión común de unidades \(m\mapsto am\), con \(a>0\),
deja invariantes ambas diferencias: el factor \(a\) se cancela en los
cocientes. La diferencia absoluta sí se multiplica por \(a\), igual que
su incertidumbre. Para un valor externo nulo, una cota unilateral o un
intervalo sin estimación central, el cociente anterior se reemplaza por
el contraste que corresponda a ese tipo de observación.

Cuando se dispone de un modelo conjunto de incertidumbre, sea
\(d_X=m_X^{\rm HMT}-m_X^{\rm ext}\). Su varianza es
\[
 u^2(d_X)=u^2(m_X^{\rm HMT})+u^2(m_X^{\rm ext})
          -2\operatorname{Cov}(m_X^{\rm HMT},m_X^{\rm ext}).
\]
La fórmula resulta de desarrollar
\(\operatorname{Var}(U-V)\). Si \(u(d_X)>0\), la diferencia
normalizada es \(z_X=d_X/u(d_X)\). Su interpretación estadística exige
el modelo probabilístico y la covarianza correspondientes; un campo no
consignado no equivale a una varianza cero. Las incertidumbres asimétricas
se conservan como tales y no se convierten automáticamente en una desviación
simétrica.

Para varias especies se mantiene la matriz de covarianza conjunta. Fijadas
las firmas discretas, si \(f_X(\vartheta)=\log(m_X/m_{\rm ref})\)
depende de coordenadas continuas ya construidas y \(J\) es su matriz
jacobiana, la propagación de primer orden da
\[
 \Sigma_f=J\Sigma_\vartheta J^{\mathsf T}.
\]
Esta fórmula conserva las correlaciones que comparten el par angular,
\(\Delta_4\), la razón de acción y las torres. Su aplicación presupone
el régimen en el que la linealización controla los términos restantes.
Un cambio entero de firma se registra como variación discreta de la
construcción, no como una fluctuación gaussiana infinitesimal.

La precisión de una evaluación tiene por ello varias fuentes: la
referencia experimental; la carta de unidades, esquema y escala; la
distribución de rutas; la aproximación de la torre; y el acoplamiento
seleccionado. Si \(\Sigma_{ij}\) denota la covarianza entre las
contribuciones \(i,j\), la covarianza total de su suma es
\(\sum_{i,j}\Sigma_{ij}\). Mantener los bloques cruzados evita
atribuir independencia a varias coincidencias que utilizan los mismos
antecedentes.

\subsection{Escala absoluta y control de la aproximación armónica}
\label{vi:subsec:vi-precision-torres}

La escala absoluta y la corrección relativa intervienen en lugares
distintos de la construcción. Fijada una sección basal de referencia
\(m_\star^{(0)}\), construida internamente, sea
\(\chi_\star=\chi_{\rm full}(\sigma_\star,\eta_\star)\).
En la normalización relativa a esa sección, la expresión es
\[
 m_X^\beta=m_\star^{(0)}\,
 \frac{\chi_{\rm full}(\sigma_X,\eta_X)}{\chi_\star}
 R_{\rm act}^{\kappa_X}.
\]
El factor basal común permanece explícito. Si la referencia es electrónica,
\(m_\star^{(0)}=m_e^{(0)}\) es la realización dimensional de
\(\mathfrak e_0\), con todos los factores de su ecuación basal, y
\(\chi_\star=\chi_{\rm full}(\sigma_e,\eta_e)\). Así se recupera
\(m_e^\beta=m_e^{(0)}R_{\rm act}^{\kappa_e}\). Si se usa el reloj
común para escribir \(m_\star^{(0)}=m_{\rm clk}b_\star\), el
factor \(b_\star\) conserva esa normalización basal; no se lo sustituye
por \(1\) ni se lo incorpora dos veces al carácter. Para dos estados
en la misma carta y con la misma referencia, el cociente es
\[
 \frac{m_Y^\beta}{m_X^\beta}
 =\frac{\chi_{\rm full}(\sigma_Y,\eta_Y)}
        {\chi_{\rm full}(\sigma_X,\eta_X)}
     R_{\rm act}^{\kappa_Y-\kappa_X}.
\]
La cancelación de \(m_\star^{(0)}/\chi_\star\) es exacta. Comprende
el prefactor basal común y los factores comunes de la rama absoluta de
acción, reloj y carta de velocidad. Por
consiguiente, la década de acción \(10^{-34}\) conserva su papel en
la realización dimensional y no se inserta otra vez en el corrector
relativo. Las diferencias entre rutas permanecen en sus firmas, lectores
y secciones de torre.

Para controlar una truncación conviene fijar la forma normal del carácter.
Sean
\[
 \zeta_{270}=135\Delta_4^2,\qquad
 N_h=\eta_h+b_{120}\mathbf 1_{\{h=120\}},\qquad
 \mathcal S=\langle90,120\rangle\setminus\{0\}.
\]
Entonces
\[
\begin{split}
 \log\chi_{\rm full}
 ={}&n_AA_{\rm rad}+\frac{s_C}{6}C^*_{\rm rad}
       +k_{\Delta_4}\Delta_4+b_\zeta\zeta_{270}
       +\sum_{h\in\mathcal S}N_hs_h .
\end{split}
\]
En esta convención, el término de altura \(120\) aparece una sola vez.
El término \(b_\zeta\zeta_{270}\) mantiene su origen cuadrático en el
registro; \(N_{270}s_{270}\) pertenece a la lectura armónica. La
igualdad del índice numérico no identifica ambos objetos.

\paragraph{Cota de truncación y error multiplicativo.}
Supóngase que, para \(h>H\), los coeficientes ya determinados de una
ruta satisfacen \(|N_hs_h|\leq Cq^h\), con \(C\geq0\) y
\(0<q<1\). La cola logarítmica \(R_H\) verifica
\[
 |R_H|\leq\sum_{\substack{h\in\mathcal S\\h>H}}Cq^h
 \leq\sum_{h=H+1}^{\infty}Cq^h
 =\frac{Cq^{H+1}}{1-q}=:\varepsilon_H.
\]
Si \(m_H\) es la coordenada obtenida al truncar esa torre y los demás
factores se mantienen fijos, \(m/m_H=e^{R_H}\). La monotonía de la
exponencial produce
\[
 e^{-\varepsilon_H}\leq\frac{m}{m_H}\leq e^{\varepsilon_H},
 \qquad
 \left|\frac{m-m_H}{m_H}\right|\leq e^{\varepsilon_H}-1.
\]
La cota permite asignar precisión a la evaluación desde la ruta y sus
operadores. Los valores de \(C,q,H\) se justifican mediante esa
construcción; el residual experimental no selecciona la longitud de la
torre ni la cota que se pretende demostrar.

% REMATE_LOCAL_VI_CONTRASTE
\Needspace{29\baselineskip}
\subsection{Conservación íntegra y reproducibilidad del contraste}
\label{vi:subsec:vi-catalogo-reproducible}

El catálogo conserva los registros completos en formatos tabulares y
estructurados. Cada fila mantiene el archivo de procedencia, su ordinal,
la clave original y el número de aparición de esa clave. Esta
identificación conserva incluso dos filas literalmente iguales: la
igualdad de sus contenidos no elimina la segunda aparición. El cambio
de formato preserva las cadenas originales, las casillas vacías y los
estados, y adjunta la huella de los bytes de entrada.

El inventario de \(855\) observables se enlaza además con las filas
metrológicas del corpus integral. La correspondencia se verifica
por el par formado por el tipo de observable y su identificador; el
control compara multiplicidades, no sólo conjuntos de claves. Las tablas
internas de celdas, familias, multisecciones y rutas permanecen separadas
de ese inventario. Las \(23\) clases, las \(20\) presentaciones
numéricas y la malla histórica conservan también sus archivos y estatutos
propios.

Los controles de transcripción comprueban la recuperación de todos los
campos, el orden de las filas, los duplicados, las cadenas decimales y la
correspondencia con las tablas de referencia. Su alcance es documental.
Las pruebas matemáticas del atlas, los lectores y la realización material
son las que establecen las salidas internas; la comparación metrológica
completa esa exposición mediante observables y cartas bien definidos.
La integridad documental y la suficiencia de una demostración se
comprueban separadamente.

La extensión del catálogo sigue, por tanto, una regla acumulativa:
una nueva evaluación incorpora el estado, su registro y firma, la sección
de torre, el lector, la realización, la carta de comparación y la
procedencia. Los casos anteriores conservan su contenido y su fecha de
corte. Así, el crecimiento del inventario hace visible una mayor parte
del dominio sin alterar retrospectivamente aquello que cada prueba y
cada evaluación habían establecido.
